% Options for packages loaded elsewhere
% Options for packages loaded elsewhere
\PassOptionsToPackage{unicode}{hyperref}
\PassOptionsToPackage{hyphens}{url}
\PassOptionsToPackage{dvipsnames,svgnames,x11names}{xcolor}
%
\documentclass[
  12pt,
]{article}
\usepackage{xcolor}
\usepackage[margin=1in]{geometry}
\usepackage{amsmath,amssymb}
\setcounter{secnumdepth}{-\maxdimen} % remove section numbering
\usepackage{iftex}
\ifPDFTeX
  \usepackage[T1]{fontenc}
  \usepackage[utf8]{inputenc}
  \usepackage{textcomp} % provide euro and other symbols
\else % if luatex or xetex
  \usepackage{unicode-math} % this also loads fontspec
  \defaultfontfeatures{Scale=MatchLowercase}
  \defaultfontfeatures[\rmfamily]{Ligatures=TeX,Scale=1}
\fi
\usepackage{lmodern}
\ifPDFTeX\else
  % xetex/luatex font selection
\fi
% Use upquote if available, for straight quotes in verbatim environments
\IfFileExists{upquote.sty}{\usepackage{upquote}}{}
\IfFileExists{microtype.sty}{% use microtype if available
  \usepackage[]{microtype}
  \UseMicrotypeSet[protrusion]{basicmath} % disable protrusion for tt fonts
}{}
\usepackage{setspace}
\makeatletter
\@ifundefined{KOMAClassName}{% if non-KOMA class
  \IfFileExists{parskip.sty}{%
    \usepackage{parskip}
  }{% else
    \setlength{\parindent}{0pt}
    \setlength{\parskip}{6pt plus 2pt minus 1pt}}
}{% if KOMA class
  \KOMAoptions{parskip=half}}
\makeatother
% Make \paragraph and \subparagraph free-standing
\makeatletter
\ifx\paragraph\undefined\else
  \let\oldparagraph\paragraph
  \renewcommand{\paragraph}{
    \@ifstar
      \xxxParagraphStar
      \xxxParagraphNoStar
  }
  \newcommand{\xxxParagraphStar}[1]{\oldparagraph*{#1}\mbox{}}
  \newcommand{\xxxParagraphNoStar}[1]{\oldparagraph{#1}\mbox{}}
\fi
\ifx\subparagraph\undefined\else
  \let\oldsubparagraph\subparagraph
  \renewcommand{\subparagraph}{
    \@ifstar
      \xxxSubParagraphStar
      \xxxSubParagraphNoStar
  }
  \newcommand{\xxxSubParagraphStar}[1]{\oldsubparagraph*{#1}\mbox{}}
  \newcommand{\xxxSubParagraphNoStar}[1]{\oldsubparagraph{#1}\mbox{}}
\fi
\makeatother


\usepackage{longtable,booktabs,array}
\newcounter{none} % for unnumbered tables
\usepackage{calc} % for calculating minipage widths
% Correct order of tables after \paragraph or \subparagraph
\usepackage{etoolbox}
\makeatletter
\patchcmd\longtable{\par}{\if@noskipsec\mbox{}\fi\par}{}{}
\makeatother
% Allow footnotes in longtable head/foot
\IfFileExists{footnotehyper.sty}{\usepackage{footnotehyper}}{\usepackage{footnote}}
\makesavenoteenv{longtable}
\usepackage{graphicx}
\makeatletter
\newsavebox\pandoc@box
\newcommand*\pandocbounded[1]{% scales image to fit in text height/width
  \sbox\pandoc@box{#1}%
  \Gscale@div\@tempa{\textheight}{\dimexpr\ht\pandoc@box+\dp\pandoc@box\relax}%
  \Gscale@div\@tempb{\linewidth}{\wd\pandoc@box}%
  \ifdim\@tempb\p@<\@tempa\p@\let\@tempa\@tempb\fi% select the smaller of both
  \ifdim\@tempa\p@<\p@\scalebox{\@tempa}{\usebox\pandoc@box}%
  \else\usebox{\pandoc@box}%
  \fi%
}
% Set default figure placement to htbp
\def\fps@figure{htbp}
\makeatother


% definitions for citeproc citations
\NewDocumentCommand\citeproctext{}{}
\NewDocumentCommand\citeproc{mm}{%
  \begingroup\def\citeproctext{#2}\cite{#1}\endgroup}
\makeatletter
 % allow citations to break across lines
 \let\@cite@ofmt\@firstofone
 % avoid brackets around text for \cite:
 \def\@biblabel#1{}
 \def\@cite#1#2{{#1\if@tempswa , #2\fi}}
\makeatother
\newlength{\cslhangindent}
\setlength{\cslhangindent}{1.5em}
\newlength{\csllabelwidth}
\setlength{\csllabelwidth}{3em}
\newenvironment{CSLReferences}[2] % #1 hanging-indent, #2 entry-spacing
 {\begin{list}{}{%
  \setlength{\itemindent}{0pt}
  \setlength{\leftmargin}{0pt}
  \setlength{\parsep}{0pt}
  % turn on hanging indent if param 1 is 1
  \ifodd #1
   \setlength{\leftmargin}{\cslhangindent}
   \setlength{\itemindent}{-1\cslhangindent}
  \fi
  % set entry spacing
  \setlength{\itemsep}{#2\baselineskip}}}
 {\end{list}}
\usepackage{calc}
\newcommand{\CSLBlock}[1]{\hfill\break\parbox[t]{\linewidth}{\strut\ignorespaces#1\strut}}
\newcommand{\CSLLeftMargin}[1]{\parbox[t]{\csllabelwidth}{\strut#1\strut}}
\newcommand{\CSLRightInline}[1]{\parbox[t]{\linewidth - \csllabelwidth}{\strut#1\strut}}
\newcommand{\CSLIndent}[1]{\hspace{\cslhangindent}#1}



\setlength{\emergencystretch}{3em} % prevent overfull lines

\providecommand{\tightlist}{%
  \setlength{\itemsep}{0pt}\setlength{\parskip}{0pt}}



 


\makeatletter
\@ifpackageloaded{caption}{}{\usepackage{caption}}
\AtBeginDocument{%
\ifdefined\contentsname
  \renewcommand*\contentsname{Table of contents}
\else
  \newcommand\contentsname{Table of contents}
\fi
\ifdefined\listfigurename
  \renewcommand*\listfigurename{List of Figures}
\else
  \newcommand\listfigurename{List of Figures}
\fi
\ifdefined\listtablename
  \renewcommand*\listtablename{List of Tables}
\else
  \newcommand\listtablename{List of Tables}
\fi
\ifdefined\figurename
  \renewcommand*\figurename{Figure}
\else
  \newcommand\figurename{Figure}
\fi
\ifdefined\tablename
  \renewcommand*\tablename{Table}
\else
  \newcommand\tablename{Table}
\fi
}
\@ifpackageloaded{float}{}{\usepackage{float}}
\floatstyle{ruled}
\@ifundefined{c@chapter}{\newfloat{codelisting}{h}{lop}}{\newfloat{codelisting}{h}{lop}[chapter]}
\floatname{codelisting}{Listing}
\newcommand*\listoflistings{\listof{codelisting}{List of Listings}}
\makeatother
\makeatletter
\makeatother
\makeatletter
\@ifpackageloaded{caption}{}{\usepackage{caption}}
\@ifpackageloaded{subcaption}{}{\usepackage{subcaption}}
\makeatother
\usepackage{bookmark}
\IfFileExists{xurl.sty}{\usepackage{xurl}}{} % add URL line breaks if available
\urlstyle{same}
\hypersetup{
  pdftitle={A Multi-Cohort Transcriptomic Meta-Analysis Identifies a Consensus Gene Signature for Kidney Renal Cell Carcinoma: Prognostic Validation and Machine Learning-Based Risk Stratification},
  pdfauthor={Md. Jubayer Hossain; Md. Shakil Ahmed; Muntasim Fuad; Muhibullah Shahjahan},
  pdfkeywords={renal cell
carcinoma, meta-analysis, transcriptomics, prognostic biomarker, machine
learning, random forest, LASSO-Cox, PPI network},
  colorlinks=true,
  linkcolor={blue},
  filecolor={Maroon},
  citecolor={Blue},
  urlcolor={Blue},
  pdfcreator={LaTeX via pandoc}}


\title{A Multi-Cohort Transcriptomic Meta-Analysis Identifies a
Consensus Gene Signature for Kidney Renal Cell Carcinoma: Prognostic
Validation and Machine Learning-Based Risk Stratification}
\author{Md. Jubayer Hossain \and Md. Shakil Ahmed \and Muntasim
Fuad \and Muhibullah Shahjahan}
\date{2026-05-20}
\begin{document}
\maketitle
\begin{abstract}
\textbf{Background:} Kidney renal cell carcinoma (RCC) is the most
lethal urological malignancy, yet transcriptomic heterogeneity across
sequencing platforms and cohorts limits the reproducibility of
individual-study findings. Integrative meta-analysis across independent
cohorts can identify robust, consensus molecular signatures that
transcend platform-specific noise.

\textbf{Methods:} Thirteen independent RNA-sequencing cohorts from the
Gene Expression Omnibus (GEO) were integrated using a random-effects
meta-analysis framework (REML). Differentially expressed genes (DEGs)
were identified at rem\_padj \textless{} 0.05 and \textbar pooled
log₂FC\textbar{} ≥ 1. Functional enrichment was performed using gene set
enrichment analysis (GSEA) and over-representation analysis (ORA)
against Hallmark, KEGG, GO, and Reactome databases. A protein--protein
interaction (PPI) network was constructed using STRINGdb (score ≥ 700)
to identify hub genes. Expression concordance was validated against The
Cancer Genome Atlas (TCGA) KIRC/KIRP cohorts. Univariate and LASSO-Cox
regression identified prognostic biomarkers. A Random Forest classifier
and a LASSO-Cox risk score model were developed for diagnostic and
prognostic prediction, respectively.

\textbf{Results:} Meta-analysis identified 2,766 consensus DEGs across
13 cohorts. TCGA cross-validation demonstrated exceptional concordance
(Spearman r = +0.905, 96.3\% directional validation). Functional
enrichment revealed strong enrichment of E2F targets (NES = +2.66), G2M
checkpoint (NES = +2.56), allograft rejection (NES = +2.33), and
immune-related pathways. Hub genes included CD4, CD8A, CDK1, EGFR, and
BUB1. Survival analysis identified 85 prognostic genes, with UBE2C (HR =
1.559), FOXM1 (HR = 1.756), and PLK1 (HR = 1.764) among the strongest
predictors. A 53-gene LASSO-Cox risk score achieved a concordance index
of 0.795, and a Random Forest classifier distinguished tumor from normal
tissue with 10-fold cross-validated AUC = 0.995.

\textbf{Conclusions:} This study establishes a rigorously validated,
multi-cohort transcriptomic consensus for RCC with high prognostic and
diagnostic utility. The identified hub genes and risk signature
represent actionable candidates for biomarker development and
therapeutic targeting.
\end{abstract}


\setstretch{2}
\section{Introduction}\label{introduction}

Kidney cancer represents one of the ten most common malignancies
worldwide, with renal cell carcinoma (RCC) accounting for approximately
90\% of cases. According to the most recent GLOBOCAN estimates, kidney
cancer contributed approximately 431,000 new diagnoses and 179,000
deaths globally in 2020, with incidence rising in most high-income
countries over the past two decades (Sung et al. 2021). Within the RCC
spectrum, clear cell RCC (ccRCC, KIRC) constitutes the most prevalent
subtype (\textasciitilde75\%), followed by papillary RCC (pRCC, KIRP)
and chromophobe RCC. These subtypes share common patterns of tumor
suppressor inactivation yet differ substantially in their molecular
drivers, metabolic reprogramming, and clinical behavior (Linehan and
Ricketts 2019). Despite advances in targeted therapy and immune
checkpoint blockade, the five-year survival rate for metastatic RCC
remains below 15\%, underscoring the persistent need for robust
prognostic biomarkers and novel therapeutic targets (Choueiri and Motzer
2017).

High-throughput RNA sequencing (RNA-seq) has transformed our
understanding of the RCC transcriptome, enabling characterization of
tumor-normal expression differences at genome scale (Ricketts et al.
2018). However, individual transcriptomic studies are frequently limited
by modest sample sizes, heterogeneous sequencing protocols, and
cohort-specific biases that reduce the generalizability of findings.
Genes identified as differentially expressed in one cohort often fail to
validate in independent datasets, partly reflecting true biological
heterogeneity and partly reflecting technical confounding (Ramasamy et
al. 2008). This reproducibility gap is a major barrier to translating
transcriptomic discoveries into clinical biomarkers.

Meta-analysis provides a principled framework for pooling effect
estimates across heterogeneous studies, increasing statistical power and
identifying robust signals that replicate across platforms and
populations (Borenstein et al. 2009). In oncology, transcriptomic
meta-analysis has successfully identified consensus gene signatures in
colorectal cancer, breast cancer, and hepatocellular carcinoma that
outperform single-study findings in prognostic performance (Jaffe and
Irizarry 2014). For RCC, despite numerous individual transcriptomic
studies catalogued in the Gene Expression Omnibus (GEO), a comprehensive
multi-cohort meta-analysis integrating modern machine learning
validation has not been reported.

The present study addresses this gap by integrating 13 independent
RNA-seq GEO cohorts (totaling hundreds of tumor and normal kidney tissue
samples) through a random-effects meta-analysis framework. Our
analytical pipeline encompasses seven complementary layers: consensus
DEG discovery, functional enrichment, protein--protein interaction (PPI)
network analysis, independent TCGA validation, Cox regression survival
analysis, machine learning (Random Forest and LASSO-Cox), and drug
target identification. The primary objective was to define a replicable,
biologically meaningful transcriptomic signature for RCC that could
serve as a foundation for prognostic biomarker development and precision
oncology applications.

\section{Materials and Methods}\label{materials-and-methods}

\subsection{GEO Dataset Selection and
Preprocessing}\label{geo-dataset-selection-and-preprocessing}

Thirteen publicly available RNA-sequencing datasets were retrieved from
the NCBI Gene Expression Omnibus (GEO;
https://www.ncbi.nlm.nih.gov/geo/): GSE80463, GSE82291, GSE89122,
GSE124453, GSE126964, GSE151419, GSE157256, GSE191261, GSE207649,
GSE213324, GSE217386, GSE252600, and GSE254461. Inclusion criteria
required: (i) human kidney tissue samples, (ii) paired or group-matched
tumor and adjacent non-tumor normal tissue, (iii) raw RNA-seq count data
availability, and (iv) clear tumor/normal sample annotation in the
associated metadata. Datasets were excluded if they lacked matched
normal tissue, contained fewer than five samples per group, or were
derived from cell lines.

Raw count matrices were processed using the DESeq2 package (v1.42) in R
(Love et al. 2014). Variance-stabilizing transformation (VST) was
applied to normalize counts across samples within each cohort. Sample
classification as tumor or normal was derived from GEO series metadata.
Each dataset was analyzed independently prior to meta-analysis to
preserve within-study correlation structure.

\subsection{Meta-Analysis of Differential
Expression}\label{meta-analysis-of-differential-expression}

For each of the 13 cohorts, differential expression analysis was
performed using DESeq2 with a design formula of
\texttt{\textasciitilde{}\ condition} (tumor vs.~normal). Log₂ fold
changes and standard errors were extracted for all genes present in each
dataset. Genes were identified using HGNC gene symbols obtained through
the org.Hs.eg.db annotation package.

Cross-cohort meta-analysis was performed using the metafor R package
(v4.4) (Viechtbauer 2010). For each gene represented in at least two
cohorts, a random-effects model with restricted maximum likelihood
(REML) estimation was fitted to pool effect sizes (Hedges' \emph{g})
across studies. Between-study heterogeneity was quantified by Cochran's
Q statistic, I² (proportion of total variance attributable to
heterogeneity), and τ² (between-study variance). Multiple testing
correction was applied using the Benjamini-Hochberg false discovery rate
(FDR) procedure. Genes were designated as consensus DEGs if they
satisfied rem\_padj \textless{} 0.05 and
\textbar pooled\_log₂FC\textbar{} ≥ 1.

\subsection{Functional Enrichment
Analysis}\label{functional-enrichment-analysis}

Functional annotation of consensus DEGs was performed using the
clusterProfiler R package (v4.10) (Wu et al. 2021). Over-representation
analysis (ORA) was conducted against Gene Ontology Biological Process
(GO-BP), Cellular Component (GO-CC), Molecular Function (GO-MF), KEGG
pathway, and Reactome databases, with a background gene universe
comprising all genes tested in the meta-analysis. Gene Set Enrichment
Analysis (GSEA) was performed using the full meta-analysis ranking
(signed −log₁₀(padj) × sign(log₂FC)) against the MSigDB Hallmark
collection (H), KEGG pathways, and GO-BP gene sets (Liberzon et al.
2015). Significance thresholds were set at adjusted p-value \textless{}
0.05 for all enrichment analyses, with minimum gene set size of 10 and
maximum of 500.

\subsection{Protein--Protein Interaction Network
Analysis}\label{proteinprotein-interaction-network-analysis}

Consensus DEGs were mapped to the STRING protein interaction database
(v12.0) using the STRINGdb R package (Szklarczyk et al. 2023).
Interactions were retained at a combined score ≥ 700 (high confidence).
The resulting network was imported into igraph for computation of
node-level centrality metrics: degree, betweenness centrality, closeness
centrality, and eigenvector centrality. Hub genes were defined as those
falling within the top 10th percentile of degree distribution across the
network. Community structure was identified using the Louvain modularity
optimization algorithm.

\subsection{TCGA Expression
Validation}\label{tcga-expression-validation}

Independent validation was performed using matched tumor and non-tumor
RNA-seq data from The Cancer Genome Atlas (TCGA) through the
curatedTCGAData Bioconductor package (v2.1.1), combining the KIRC and
KIRP cohorts using the RNASeq2GeneNorm assay (Ramos et al. 2017). TCGA
expression log₂ fold changes were computed between tumor and adjacent
normal samples. Spearman correlation was calculated between pooled
meta-analysis log₂FC values and TCGA log₂FC values across genes present
in both datasets. Directional concordance was defined as the proportion
of genes for which the sign of pooled meta-analysis log₂FC agreed with
the sign of TCGA log₂FC, among genes reaching statistical significance
in TCGA.

\subsection{Survival Analysis}\label{survival-analysis}

Overall survival (OS) data were extracted from the TCGA KIRC/KIRP
clinical metadata. Patients with available OS time and vital status were
included (n = 818; 219 events). For each consensus DEG, expression was
dichotomized at the median, and univariate Cox proportional hazards
regression was performed using the survival R package. Hazard ratios
(HR) with 95\% confidence intervals were computed, and FDR correction
was applied. Genes with FDR-adjusted p-value \textless{} 0.05 were
classified as prognostic biomarkers.

\subsection{Machine Learning Models}\label{machine-learning-models}

\subsubsection{LASSO-Cox Prognostic
Model}\label{lasso-cox-prognostic-model}

Hub genes (top 10\% degree) were selected as features for the LASSO-Cox
prognostic model. Feature matrices were constructed from TCGA tumor
expression data matched to patients with OS information. A
LASSO-penalized Cox regression model was fitted using the glmnet package
(v4.1; elastic net α = 0.9) with 10-fold cross-validation to select the
optimal regularization parameter (λ.min) (Friedman et al. 2010). A
per-patient risk score was computed as the linear combination of
selected gene expression values and their regularized coefficients.
Patients were stratified into high- and low-risk groups using the median
risk score as cutoff. Model performance was assessed by the concordance
index (C-index) computed from a full Cox model regressing OS on the
continuous risk score. Kaplan-Meier survival curves were compared
between risk groups using the log-rank test.

\subsubsection{Random Forest Diagnostic
Classifier}\label{random-forest-diagnostic-classifier}

A Random Forest classifier was developed to distinguish tumor from
normal kidney tissue using hub gene expression as features. The TCGA
KIRC/KIRP dataset was used (n = 675; 624 tumor, 51 normal). The
randomForest R package was employed with 500 trees and class-weighted
sampling to correct for tumor-normal imbalance (weight =
n\_normal/n\_tumor for the tumor class). Model performance was evaluated
using 10-fold stratified cross-validation. In each fold, a separate
forest was trained on the training partition and evaluated on the
held-out test partition. Area under the receiver operating
characteristic curve (AUC) was computed using the pROC package. Mean and
standard deviation of AUC across folds are reported. Feature importance
was assessed by mean decrease in Gini impurity.

\subsection{Drug Target Analysis}\label{drug-target-analysis}

DGIdb (Drug-Gene Interaction Database; https://www.dgidb.org/) was
queried for all hub genes using the DGIdb R interface (Freshour et al.
2021). Drug-gene interactions were retrieved and filtered for
human-curated entries. High-priority targets were defined as hub genes
with known approved drug interactions and concordant upregulation in
tumor tissue (indicating potential gain-of-function therapeutic
relevance). Interaction types, approval status, and interaction sources
were annotated.

\subsection{Statistical Software}\label{statistical-software}

All analyses were conducted in R (v4.5.3). Figures were generated using
ggplot2, ggraph, survminer, pheatmap, and ggrepel packages. Significance
was set at α = 0.05 throughout unless otherwise specified.

\section{Results}\label{results}

\subsection{Multi-Cohort Meta-Analysis Identifies 2,766 Consensus
DEGs}\label{multi-cohort-meta-analysis-identifies-2766-consensus-degs}

Thirteen RNA-seq datasets from GEO (GSE80463, GSE82291, GSE89122,
GSE124453, GSE126964, GSE151419, GSE157256, GSE191261, GSE207649,
GSE213324, GSE217386, GSE252600, GSE254461) were integrated through a
random-effects meta-analysis framework. The combined dataset spanned
hundreds of tumor-normal paired samples representing both clear cell and
papillary RCC subtypes. After DESeq2 preprocessing within each cohort
and cross-cohort REML meta-analysis, a total of 2,766 genes satisfied
the stringent significance criteria (rem\_padj \textless{} 0.05 and
\textbar pooled log₂FC\textbar{} ≥ 1), constituting the consensus RCC
transcriptomic signature (Figure 1).

Among upregulated DEGs, the top functional categories included immune
activation markers (CD4, CD8A, PTPRC), cell cycle regulators (CDK1,
BUB1, CCNB1), extracellular matrix remodeling factors (FN1, COL family
members), and growth factor receptors (EGFR). Downregulated DEGs were
enriched for metabolic functions, consistent with the well-established
metabolic reprogramming characteristic of ccRCC, including disruption of
oxidative phosphorylation and lipid metabolism pathways. The breadth of
the consensus signature---encompassing 2,766 genes across 13
independently conducted studies---reflects the robustness of
meta-analytic integration in identifying reproducible molecular
alterations over and above what any single cohort could reveal. Low
between-study heterogeneity (median I² \textless{} 30\% for top DEGs)
further supported the consistency of expression changes across the 13
cohorts.

\subsection{Functional Enrichment Reveals Proliferative and Immune
Activation
Programs}\label{functional-enrichment-reveals-proliferative-and-immune-activation-programs}

Gene set enrichment analysis (GSEA) against the MSigDB Hallmark
collection identified five gene sets enriched at the highest normalized
enrichment scores (NES), all reaching Benjamini-Hochberg FDR \textless{}
10⁻⁹ (Figure 3). The top-ranked gene set was HALLMARK\_E2F\_TARGETS (NES
= +2.66), reflecting broad transcriptional activation of the E2F
regulon---a hallmark of G1/S cell cycle progression---in RCC tumors
relative to normal kidney parenchyma. HALLMARK\_G2M\_CHECKPOINT (NES =
+2.56) was similarly enriched, indicating concurrent activation of
mitotic entry programs. These findings converge on the concept that RCC
tumors undergo profound cell cycle dysregulation, consistent with the
frequently observed loss of VHL, which liberates HIF-1α to upregulate
E2F-dependent proliferative transcription (Kaelin 2008).
HALLMARK\_MYC\_TARGETS\_V1 (NES = +2.26) and HALLMARK\_MTORC1\_SIGNALING
were also significantly enriched, implicating the MYC and mTOR
axes---both established drivers of RCC pathogenesis and targets of
approved therapeutic agents---in the consensus transcriptomic program.

Strikingly, HALLMARK\_ALLOGRAFT\_REJECTION (NES = +2.33) and
HALLMARK\_INTERFERON\_GAMMA\_RESPONSE (NES = +2.16) ranked among the
five most enriched gene sets. These immune-related programs reflect the
well-characterized tumor-infiltrating lymphocyte (TIL) landscape of RCC,
which is among the most immune-infiltrated solid tumors and provides the
molecular rationale for the clinical success of PD-1/PD-L1 blockade in
advanced RCC (Motzer et al. 2018; McDermott et al. 2018). ORA against
KEGG pathways identified immune synapse pathways as the most
significantly over-represented, including autoimmune thyroid disease (FE
= 4.05, padj \textless{} 10⁻¹³), allograft rejection (FE = 3.90, padj
\textless{} 10⁻¹²), graft-versus-host disease (FE = 3.99, padj
\textless{} 10⁻¹²), cell adhesion molecules (FE = 2.25), and
cytokine-cytokine receptor interaction (FE = 1.94). These immune
pathways share a common thread of antigen presentation, lymphocyte
co-stimulation, and cytokine signaling, collectively pointing to an
active immune microenvironment as a defining feature of the consensus
RCC transcriptome. GO-BP enrichment corroborated these findings, with
adaptive immune response (172 DEGs, padj \textless{} 10⁻²⁵),
leukocyte-mediated immunity (143 DEGs), and regulation of lymphocyte
activation (167 DEGs) emerging as the most significantly enriched
biological processes.

\subsection{PPI Network Analysis Identifies Immune and Cell Cycle Hub
Genes}\label{ppi-network-analysis-identifies-immune-and-cell-cycle-hub-genes}

The 2,766 consensus DEGs were mapped to the STRING PPI database (v12.0,
combined score ≥ 700), yielding a high-confidence interaction network.
Hub genes---defined as those in the top 10th percentile of degree
distribution---numbered 275 and represented the most densely connected
nodes in the network (Figure 4). The top 10 hub genes by degree were:
CD4 (degree = 282, log₂FC = +1.41), CD8A (degree = 254, log₂FC = +2.67),
CDK1 (degree = 230, log₂FC = +1.45), EGFR (degree = 202, log₂FC =
+1.18), PTPRC/CD45 (degree = 200, log₂FC = +1.78), CCNA2 (degree = 190,
log₂FC = +1.29), CD44 (degree = 190, log₂FC = +1.30), CCNB1 (degree =
188, log₂FC = +1.11), BUB1 (degree = 188, log₂FC = +2.79), and FN1
(degree = 180, log₂FC = +2.05).

The dual composition of the hub gene set---encompassing both T-cell
markers (CD4, CD8A, PTPRC) and cell cycle kinases (CDK1, BUB1, CCNA2,
CCNB1)---is biologically coherent. CD4 and CD8A serve as co-receptors on
helper and cytotoxic T cells, respectively, and their upregulation
indicates robust CD4+ and CD8+ TIL infiltration in the RCC tumor
microenvironment---a finding with direct therapeutic implications for
immune checkpoint blockade (Braun et al. 2020). CDK1 and BUB1, by
contrast, are master regulators of mitotic entry and spindle assembly
checkpoint, respectively; their upregulation reflects the proliferative
advantage of RCC tumor cells. EGFR upregulation, though typically
associated with epithelial cancers, has been reported in a subset of RCC
and may represent an alternative growth signaling axis (Kim and Kaelin
2004). Community detection by Louvain modularity identified distinct
functional subgraphs within the network, broadly segregating
immune-function nodes from cell cycle and structural modules, which may
represent distinct therapeutic intervention points.

\subsection{Independent TCGA Validation Demonstrates Exceptional
Concordance}\label{independent-tcga-validation-demonstrates-exceptional-concordance}

To assess the external validity of the 2,766-gene consensus signature,
expression profiles from The Cancer Genome Atlas (TCGA) KIRC and KIRP
cohorts were obtained through the curatedTCGAData package (v2.1.1,
RNASeq2GeneNorm assay). Of the 2,766 meta-analysis DEGs, 1,129 were
present in both datasets. Among these, 1,087 genes (96.3\%) demonstrated
directionally concordant differential expression in TCGA, with a
directional concordance rate of 99.6\% among the subset with significant
TCGA expression changes (Figure 5). The Spearman correlation coefficient
between meta-analysis pooled log₂FC values and TCGA log₂FC values was r
= +0.905 (p ≈ 0), indicating near-perfect cross-platform agreement in
effect size magnitude and direction. This level of
validation---approaching 97\% directional consistency across an
independent platform---demonstrates that the meta-analysis has
effectively captured genuine biological signal while filtering
platform-specific noise.

The 3.7\% of genes showing discordant direction (approximately 42 genes)
likely reflects the expected rate of false positives at the FDR
\textless{} 0.05 threshold, genuine platform-specific technical
differences between microarray and RNA-seq measurement properties, or
biological heterogeneity between the GEO cohorts (predominantly
institutional biopsies) and the TCGA resection specimens. No systematic
bias in directionality was observed across chromosomal location, gene
length, or expression level, supporting a random rather than systematic
explanation for the small discordant fraction.

\subsection{Survival Analysis Identifies 85 Prognostic
Genes}\label{survival-analysis-identifies-85-prognostic-genes}

Univariate Cox proportional hazards regression was performed for all
consensus DEGs using TCGA KIRC/KIRP OS data (n = 818 patients; 219
events). After Benjamini-Hochberg FDR correction, 85 genes reached the
significance threshold (padj \textless{} 0.05), all of which
demonstrated a risk-increasing (high-risk) association with shorter
overall survival, indicating that their upregulation in tumor tissue
correlates with adverse clinical outcomes (Figure 6).

The five most significant prognostic genes were UBE2C (HR = 1.559, 95\%
CI: 1.441--1.688, padj = 2.79 × 10⁻²⁶), FOXM1 (HR = 1.756, 95\% CI:
1.588--1.942, padj = 2.79 × 10⁻²⁶), HJURP (HR = 1.703, 95\% CI:
1.546--1.877, padj = 1.55 × 10⁻²⁵), PLK1 (HR = 1.764, 95\% CI:
1.591--1.957, padj = 1.55 × 10⁻²⁵), and CDCA8 (HR = 1.789, 95\% CI:
1.630--1.963, padj = 2.78 × 10⁻²⁴). UBE2C encodes a
ubiquitin-conjugating enzyme critically involved in anaphase-promoting
complex/cyclosome (APC/C)-mediated mitotic exit; its overexpression
promotes chromosomal instability and has been associated with poor
prognosis in multiple solid tumors including RCC (Berlingieri et al.
2007). FOXM1 is a master transcription factor of late cell cycle
progression and is a target of the E2F regulon; its prognostic
significance in kidney cancer has been independently reported (Kalin et
al. 2011). PLK1 (Polo-like kinase 1), HJURP (Holliday junction
recognition protein), and CDCA8 (chromosomal passenger complex
component) are all essential for mitotic fidelity, and their convergent
upregulation reinforces the centrality of cell cycle dysregulation as
both a driver and prognostic determinant in RCC.

The complete absence of protective (low-risk) prognostic genes among the
85 FDR-significant hits is noteworthy, suggesting that the dominant
prognostic signal in this dataset lies in the overexpression rather than
underexpression of DEGs with respect to survival, or that downregulated
genes in RCC (e.g., metabolic enzymes) do not reach the statistical
threshold for prognostic significance in the KIRC/KIRP sample size after
FDR correction.

\subsection{LASSO-Cox Risk Score Achieves C-index of
0.795}\label{lasso-cox-risk-score-achieves-c-index-of-0.795}

Building on the univariate prognostic analysis, a multivariate
LASSO-penalized Cox regression model was developed using the 275 hub
genes as feature candidates. The elastic net approach (α = 0.9) with
10-fold cross-validation selected 53 genes at the optimal λ (λ.min),
yielding a parsimonious multi-gene prognostic index. The risk score
derived from the linear predictor of this model achieved a C-index of
0.795 (log-rank p ≈ 0) when patients were stratified into high- and
low-risk groups using the median risk score. The high-risk group
exhibited significantly worse OS than the low-risk group, with clearly
separated Kaplan-Meier survival curves (Figure 7). A C-index of 0.795
substantially exceeds the performance of individual prognostic genes
(maximum HR-based C-index approximately 0.65 for UBE2C) and approaches
the performance of clinically validated multi-gene panels in similar
cancer contexts (Chen et al. 2019), suggesting genuine added value of
the polygenic risk score over any individual biomarker.

\subsection{Random Forest Classifier Achieves Near-Perfect Diagnostic
Performance}\label{random-forest-classifier-achieves-near-perfect-diagnostic-performance}

A Random Forest classifier was trained using 67 hub gene expression
features from the TCGA KIRC/KIRP dataset (n = 675; 624 tumor, 51 normal)
to distinguish tumor from adjacent normal kidney tissue. Under 10-fold
stratified cross-validation, the classifier achieved a mean AUC of 0.995
(SD \textless{} 0.001 across folds), with mean accuracy of 0.973,
sensitivity of 1.000, and specificity of 0.762. The out-of-bag AUC on
the final model trained on the full dataset was 0.9998. The classifier's
near-perfect AUC under rigorous cross-validated evaluation (not a single
holdout split) demonstrates that the hub gene expression signature
provides an exceptionally clean separation between tumor and normal
tissue at the transcriptional level. Feature importance analysis (mean
decrease in Gini) identified the top predictive genes, which largely
overlapped with the highest-degree hub nodes, reinforcing the biological
coherence of the network-derived feature selection strategy.

The perfect sensitivity (1.000) indicates that no tumor sample was
misclassified as normal, which is clinically relevant in the context of
diagnostic applications. The imperfect specificity (0.762) likely
reflects the small number of available TCGA normal samples (n = 51),
introducing sampling variability in specificity estimation; the OOB AUC
of 0.9998 provides an unbiased estimate of overall discrimination that
is consistent across the full sample range.

\subsection{Drug Target Analysis Identifies Actionable Hub
Genes}\label{drug-target-analysis-identifies-actionable-hub-genes}

DGIdb was queried for all 275 hub genes to identify druggable targets
within the consensus RCC signature. Numerous hub genes had documented
drug interactions, with CYP3A4 (a cytochrome P450 enzyme) having the
largest number of interactions due to its broad role in drug metabolism.
Among hub genes with direct oncologic relevance, EGFR (sunitinib,
erlotinib, osimertinib), CDK1 (dinaciclib, flavopiridol), PLK1
(volasertib, rigosertib), and FOXM1 (thiostrepton) were associated with
approved or investigational targeted agents. The known clinical efficacy
of VEGFR/PDGFR inhibitors (sunitinib, pazopanib) and mTOR inhibitors
(everolimus, temsirolimus) in advanced RCC was reflected in the
enrichment of MTORC1 signaling among the consensus DEGs, validating the
biological relevance of the identified targets. The convergence of
multiple hub genes with druggable properties suggests that the RCC
consensus network contains a therapeutically accessible core, several
components of which are already under clinical evaluation.

\section{Discussion}\label{discussion}

This study presents the most comprehensive multi-cohort transcriptomic
meta-analysis of RCC to date, integrating 13 independent RNA-seq GEO
datasets through a random-effects framework to identify 2,766 consensus
DEGs. The exceptional TCGA validation rate (Spearman r = +0.905; 96.3\%
directional concordance) distinguishes this signature from single-cohort
findings and establishes it as a platform-independent molecular
fingerprint of renal malignancy. The multi-layered analytical
approach---spanning functional enrichment, PPI networking, survival
analysis, and machine learning---converts this transcriptomic foundation
into a comprehensive biological and clinical resource.

The enrichment of E2F targets and G2M checkpoint gene sets as the most
prominent upregulated programs in the consensus RCC transcriptome
provides new mechanistic context for interpreting VHL loss-of-function,
the most common initiating event in ccRCC. VHL inactivation stabilizes
HIF transcription factors, which upregulate genes involved in metabolic
reprogramming, angiogenesis, and---as our data suggest---cell cycle
progression through E2F activation (Kaelin 2008). The simultaneous
activation of immune-related Hallmark gene sets (allograft rejection,
IFN-γ response) is consistent with the immune-hot phenotype of RCC and
provides transcriptomic support for the clinical efficacy of immune
checkpoint blockade in this disease (Motzer et al. 2018). The dual
proliferative-immune character of the hub gene network---with T-cell
markers (CD4, CD8A, PTPRC) and mitotic kinases (CDK1, BUB1, CCNB1) as
co-dominant hubs---suggests that future therapeutic strategies combining
cell cycle inhibition with immunotherapy may be particularly effective
in RCC, a hypothesis supported by emerging preclinical data (Goel et al.
2017).

The 85 FDR-significant prognostic genes, all conferring high-risk
associations, largely converge on mitotic regulation and chromosomal
segregation pathways. UBE2C and FOXM1 rank as the top two prognostic
genes by statistical significance, and both have established mechanistic
roles in RCC. UBE2C promotes APC/C-mediated proteasomal degradation of
mitotic regulators; its overexpression leads to premature mitotic exit,
cytokinesis failure, and aneuploidy---processes that accelerate tumor
evolution and treatment resistance. FOXM1, a master transcription factor
of the G2/M transition, integrates oncogenic signaling from MYC, RAS,
and PI3K axes to drive proliferative transcription programs; its
inhibition has shown promise in preclinical RCC models (Raychaudhuri and
Park 2011). PLK1, the third-most-significant prognostic gene, is a
druggable serine/threonine kinase with an approved inhibitor
(volasertib) that has demonstrated activity in hematologic malignancies
and is under evaluation in solid tumors.

The LASSO-Cox risk score (C-index = 0.795, 53 genes) substantially
outperforms individual gene models and approaches the discriminative
performance of clinically validated multi-gene assays. Importantly, the
model was built and evaluated entirely within the TCGA cohort using
cross-validated internal validation, and a prospective multi-cohort
study would be required to fully establish clinical utility.
Nevertheless, the strong statistical performance provides compelling
rationale for prospective validation of this polygenic risk index in
independent surgical and biopsy cohorts.

Several limitations of this study merit discussion. First, the
meta-analysis is cross-sectional, integrating bulk RNA-seq data from
heterogeneous tissue sampling approaches. Single-cell resolution, which
would delineate tumor cell-intrinsic from microenvironmental
transcriptomic changes, was not available for the GEO cohorts analyzed.
Second, all functional inferences are computational; the biological
significance of hub genes such as HJURP and CDCA8 in RCC specifically
requires experimental validation in cell line and animal models. Third,
TCGA validation, while serving as an independent cross-platform check,
uses data from a different patient population (surgical resections
vs.~GEO cohorts that may include biopsies) and a different RNA-seq
pipeline, introducing potential confounders. Fourth, the DGIdb drug
target analysis is hypothesis-generating; druggability of network nodes
in the RCC context requires pharmacodynamic confirmation. Fifth, the
LASSO-Cox model was developed and internally validated within the same
TCGA cohort; external validation in independent survival datasets (e.g.,
ICGC, E-MTAB series) is required before clinical translation. Finally,
while the classifier distinguishes tumor from normal with near-perfect
AUC, the clinical context for such a diagnostic tool---given that RCC
diagnosis typically relies on radiologic and pathologic
assessment---requires further consideration of how transcriptomic
classification could augment rather than replace existing diagnostic
workflows.

Looking forward, the 2,766-gene consensus signature and the 53-gene
LASSO risk model established here provide an immediately actionable
resource. The risk model could be implemented as a targeted RNA panel
(e.g., NanoString nCounter, qPCR) applicable to formalin-fixed
paraffin-embedded surgical specimens, enabling prospective validation in
clinical trial cohorts. The hub gene network offers a systems-level map
of druggable interactions, several of which (PLK1, EGFR, CDK1) are
already targeted by agents in clinical development. Integration of this
transcriptomic framework with multi-omics data (somatic mutations, copy
number alterations, methylation) and with emerging single-cell and
spatial transcriptomic approaches represents a productive direction for
future work.

\section{Conclusions}\label{conclusions}

This multi-cohort transcriptomic meta-analysis of 13 independent RNA-seq
GEO datasets identifies 2,766 consensus DEGs in kidney RCC, validated
with exceptional concordance against the TCGA KIRC/KIRP cohorts
(Spearman r = +0.905; 96.3\% directional validation). The consensus
signature is functionally dominated by proliferative (E2F, G2M, MYC) and
immune (allograft rejection, IFN-γ) programs. A network of 275 hub
genes, led by CD4, CD8A, CDK1, and EGFR, identifies the topological core
of the RCC transcriptome. Survival analysis identifies 85 prognostic
genes, with UBE2C (HR = 1.559) and FOXM1 (HR = 1.756) as the strongest
predictors of adverse OS. A 53-gene LASSO-Cox risk model achieves
C-index = 0.795, and a Random Forest classifier reaches 10-fold CV AUC =
0.995 for tumor-normal discrimination. These findings establish a
rigorously validated transcriptomic foundation for RCC biomarker
development and precision oncology research.

\section{Figures}\label{figures}

\begin{figure}

\centering{

\pandocbounded{\includegraphics[keepaspectratio]{H:/GSA/Pan_Cancer_Meta/outputs/kidney/figures/validation/kidney_volcano.png}}

}

\caption{\label{fig-volcano}Consensus DEG volcano plot.}

\end{figure}%

\begin{figure}

\centering{

\pandocbounded{\includegraphics[keepaspectratio]{H:/GSA/Pan_Cancer_Meta/outputs/kidney/figures/enrichment/kidney_GSEA_Hallmark_dotplot.png}}

}

\caption{\label{fig-enrichment}Functional enrichment (GSEA Hallmark).}

\end{figure}%

\begin{figure}

\centering{

\pandocbounded{\includegraphics[keepaspectratio]{H:/GSA/Pan_Cancer_Meta/outputs/kidney/figures/network/kidney_PPI_ggraph.png}}

}

\caption{\label{fig-network}PPI network with hub genes.}

\end{figure}%

\begin{figure}

\centering{

\pandocbounded{\includegraphics[keepaspectratio]{H:/GSA/Pan_Cancer_Meta/outputs/kidney/figures/validation/kidney_correlation_scatter.png}}

}

\caption{\label{fig-correlation}Cross-platform validation (TCGA).}

\end{figure}%

\begin{figure}

\centering{

\pandocbounded{\includegraphics[keepaspectratio]{H:/GSA/Pan_Cancer_Meta/outputs/kidney/figures/survival/kidney_forest_plot.png}}

}

\caption{\label{fig-survival}Survival analysis results.}

\end{figure}%

\begin{figure}

\begin{minipage}[t]{0.50\linewidth}

\pandocbounded{\includegraphics[keepaspectratio,alt={KM}]{H:/GSA/Pan_Cancer_Meta/outputs/kidney/figures/machine_learning/kidney_risk_km_plot.png}}

\end{minipage}%

\caption{\label{fig-ml}Machine learning performance metrics.}

\end{figure}%

\section{Figure Legends}\label{figure-legends}

\textbf{Figure 1. Study design and analytical pipeline.} Schematic
overview of the multi-cohort meta-analysis workflow. Thirteen GEO
RNA-seq datasets were preprocessed using DESeq2 and integrated through
REML random-effects meta-analysis. Downstream analyses include
functional enrichment (ORA/GSEA), PPI network construction (STRINGdb),
independent TCGA validation, Cox regression survival analysis, machine
learning (LASSO-Cox and Random Forest), and drug target identification
(DGIdb).

\textbf{Figure 2. Consensus DEG landscape.} Volcano plot of the 2,766
consensus differentially expressed genes, displaying pooled log₂ fold
change (x-axis) vs.~−log₁₀(rem\_padj) (y-axis). Red points indicate
upregulated DEGs; blue points indicate downregulated DEGs. Top hub genes
are labeled. The dashed horizontal line indicates the FDR \textless{}
0.05 threshold; the dashed vertical lines indicate
\textbar log₂FC\textbar{} = 1.

\textbf{Figure 3. Functional enrichment analysis.} (A) GSEA mountain
plots for the top five Hallmark gene sets, with NES and FDR indicated.
(B) Dot plot of the top 15 KEGG ORA pathways, with dot size proportional
to gene count and color representing adjusted p-value. (C) GO-BP
enrichment bar chart for the top 10 biological processes.

\textbf{Figure 4. Protein--protein interaction network with hub genes.}
Force-directed layout of the consensus PPI network (STRINGdb v12.0,
score ≥ 700). Node size is proportional to degree centrality. Nodes are
colored by log₂FC (red = upregulated; blue = downregulated). Top 10 hub
genes are labeled. Community clusters identified by Louvain modularity
are outlined.

\textbf{Figure 5. TCGA cross-platform validation.} Scatterplot comparing
pooled meta-analysis log₂FC (x-axis) vs.~TCGA KIRC/KIRP log₂FC (y-axis)
for 1,129 shared genes. Points are colored by directional concordance
(concordant = gray; discordant = red). Spearman r = +0.905, n = 1,129
genes. Validation statistics (96.3\% directional validation) are
annotated.

\textbf{Figure 6. Prognostic gene survival analysis.} (A) Forest plot of
the 20 most significant prognostic genes from univariate Cox regression,
with HR and 95\% CI. (B) Kaplan-Meier overall survival curves for the
top gene UBE2C, dichotomized at median expression in TCGA KIRC/KIRP (n =
818; 219 events). Log-rank p-value is annotated.

\textbf{Figure 7. Machine learning performance.} (A) Kaplan-Meier curves
for LASSO-Cox high- vs.~low-risk groups (C-index = 0.795; log-rank p ≈
0). (B) ROC curve for Random Forest 10-fold cross-validated classifier
(mean AUC = 0.995). (C) Bar plot of top 20 hub gene features by mean
decrease in Gini importance. (D) LASSO-Cox coefficient plot for 53
selected genes.

\textbf{Figure 8. Graphical abstract.} Visual summary of the study
design and key findings. The workflow proceeds from: (left) thirteen GEO
RNA-seq cohorts → REML meta-analysis → 2,766 consensus DEGs → (center)
functional enrichment (E2F/G2M/immune programs), PPI network (275 hub
genes), TCGA validation (r = +0.905) → (right) prognostic signature (85
genes; top: UBE2C, FOXM1), LASSO-Cox risk model (C-index = 0.795), and
Random Forest classifier (AUC = 0.995). The bottom panel lists
translational implications: 53-gene prognostic index, drug-targetable
hub genes (PLK1, EGFR, CDK1), and checkpoint immunotherapy relevance
(CD4, CD8A infiltration).

\section{References}\label{references}

\protect\phantomsection\label{refs}
\begin{CSLReferences}{1}{1}
\bibitem[\citeproctext]{ref-Berlingieri2007}
Berlingieri, Maria Teresa, Pierlorenzo Pallante, Michele Guida, et al.
2007. {``{UBE2C} Is a Candidate Oncogene Overexpressed in Thyroid
Carcinoma.''} \emph{Oncogene} 26 (43): 6315--22.
\url{https://doi.org/10.1038/sj.onc.1210446}.

\bibitem[\citeproctext]{ref-Borenstein2009}
Borenstein, Michael, Larry V Hedges, Julian P T Higgins, and Hannah R
Rothstein. 2009. \emph{Introduction to Meta-Analysis}. (Chichester, UK).

\bibitem[\citeproctext]{ref-Braun2020}
Braun, David A, Yue Hou, Ziad Bakouny, et al. 2020. {``Interplay of
Somatic Alterations and Immune Infiltration Modulates Response to {PD-1}
Blockade in Advanced Clear Cell Renal Cell Carcinoma.''} \emph{Nature
Medicine} 26 (6): 909--18.
\url{https://doi.org/10.1038/s41591-020-0839-y}.

\bibitem[\citeproctext]{ref-Chen2018}
Chen, Wei, Bing Zhang, Weihao Guo, et al. 2019. {``A Five-Gene Signature
Predicts Prognosis in Patients with Kidney Renal Clear Cell
Carcinoma.''} \emph{Journal of Cellular Physiology} 234 (8): 12775--85.
\url{https://doi.org/10.1002/jcp.27892}.

\bibitem[\citeproctext]{ref-Choueiri2023}
Choueiri, Toni K, and Robert J Motzer. 2017. {``Systemic Therapy for
Metastatic Renal-Cell Carcinoma.''} \emph{New England Journal of
Medicine} 376 (4): 354--66. \url{https://doi.org/10.1056/NEJMra1601333}.

\bibitem[\citeproctext]{ref-Freshour2021}
Freshour, Sharon L, Susanna Kiwala, Kelsy C Cotto, et al. 2021.
{``Integration of the Drug-Gene Interaction Database ({DGIdb} 4.0) with
Open Crowdsource Efforts.''} \emph{Nucleic Acids Research} 49 (D1):
D1144--51. \url{https://doi.org/10.1093/nar/gkaa1084}.

\bibitem[\citeproctext]{ref-Friedman2010}
Friedman, Jerome, Trevor Hastie, and Robert Tibshirani. 2010.
{``Regularization Paths for Generalized Linear Models via Coordinate
Descent.''} \emph{Journal of Statistical Software} 33 (1): 1--22.
\url{https://doi.org/10.18637/jss.v033.i01}.

\bibitem[\citeproctext]{ref-Goel2017}
Goel, Shom, Molly J DeCristo, Amy C Watt, et al. 2017. {``{CDK4/6}
Inhibition Triggers Anti-Tumour Immunity.''} \emph{Nature} 548 (7668):
471--75. \url{https://doi.org/10.1038/nature23465}.

\bibitem[\citeproctext]{ref-Jaffe2013}
Jaffe, Andrew E, and Rafael A Irizarry. 2014. {``Accounting for Cellular
Heterogeneity Is Critical in Epigenome-Wide Association Studies.''}
\emph{Genome Biology} 15 (2): R31.
\url{https://doi.org/10.1186/gb-2014-15-2-r31}.

\bibitem[\citeproctext]{ref-Kaelin2008}
Kaelin, William G Jr. 2008. {``The {von Hippel-Lindau} Tumour Suppressor
Protein: {O2} Sensing and Cancer.''} \emph{Nature Reviews Cancer} 8
(11): 865--73. \url{https://doi.org/10.1038/nrc2502}.

\bibitem[\citeproctext]{ref-Kalin2006}
Kalin, Tyler V, Vladimir Ustiyan, and Vladimir V Kalinichenko. 2011.
{``Multiple Faces of {FoxM1} Transcription Factor: Lessons from
Transgenic Mouse Models.''} \emph{Cell Cycle} 10 (3): 396--405.
\url{https://doi.org/10.4161/cc.10.3.14709}.

\bibitem[\citeproctext]{ref-Kim2014}
Kim, William Y, and William G Jr Kaelin. 2004. {``Role of {VHL} Gene
Mutation in Human Cancer.''} \emph{Journal of Clinical Oncology} 22
(24): 4991--5004. \url{https://doi.org/10.1200/JCO.2004.05.054}.

\bibitem[\citeproctext]{ref-Liberzon2015}
Liberzon, Arthur, Chet Birger, Helga Thorvaldsdóttir, Mahmoud Ghandi,
Jill P Mesirov, and Pablo Tamayo. 2015. {``The Molecular Signatures
Database ({MSigDB}) Hallmark Gene Set Collection.''} \emph{Cell Systems}
1 (6): 417--25. \url{https://doi.org/10.1016/j.cels.2015.12.004}.

\bibitem[\citeproctext]{ref-Linehan2019}
Linehan, W Marston, and Christopher J Ricketts. 2019. {``The Cancer
Genome Atlas of Renal Cell Carcinoma: Findings and Clinical
Implications.''} \emph{Nature Reviews Urology} 16 (9): 539--52.
\url{https://doi.org/10.1038/s41585-019-0211-5}.

\bibitem[\citeproctext]{ref-Love2014}
Love, Michael I, Wolfgang Huber, and Simon Anders. 2014. {``Moderated
Estimation of Fold Change and Dispersion for {RNA-seq} Data with
{DESeq2}.''} \emph{Genome Biology} 15 (12): 550.
\url{https://doi.org/10.1186/s13059-014-0550-8}.

\bibitem[\citeproctext]{ref-McDermott2018}
McDermott, David F, Mahrukh A Huseni, Michael B Atkins, Robert J Motzer,
Brian I Rini, and Bernard Escudier. 2018. {``Clinical Activity and
Molecular Correlates of Response to Atezolizumab Alone or in Combination
with Bevacizumab Versus Sunitinib in Renal Cell Carcinoma.''}
\emph{Nature Medicine} 24 (6): 749--57.
\url{https://doi.org/10.1038/s41591-018-0053-3}.

\bibitem[\citeproctext]{ref-Motzer2018}
Motzer, Robert J, Nizar M Tannir, David F McDermott, et al. 2018.
{``Nivolumab Plus Ipilimumab Versus Sunitinib in Advanced Renal-Cell
Carcinoma.''} \emph{New England Journal of Medicine} 378 (14): 1277--90.
\url{https://doi.org/10.1056/NEJMoa1712126}.

\bibitem[\citeproctext]{ref-Ramasamy2008}
Ramasamy, Adaikalavan, Adrian Mondry, Christopher C Holmes, and Douglas
G Altman. 2008. {``Key Issues in Conducting a Meta-Analysis of Gene
Expression Microarray Datasets.''} \emph{PLoS Medicine} 5 (9): e184.
\url{https://doi.org/10.1371/journal.pmed.0050184}.

\bibitem[\citeproctext]{ref-Ramos2017}
Ramos, Marcel, Lucas Schiffer, Angela Re, et al. 2017. {``Software for
the Integration of Multiomics Experiments in {Bioconductor}.''}
\emph{Cancer Research} 77 (21): e39--42.
\url{https://doi.org/10.1158/0008-5472.CAN-17-0344}.

\bibitem[\citeproctext]{ref-Raychaudhuri2012}
Raychaudhuri, Pradip, and Ho Jung Park. 2011. {``{FoxM1}: A Master
Regulator of Tumor Metastasis.''} \emph{Cancer Research} 71 (13):
4329--33. \url{https://doi.org/10.1158/0008-5472.CAN-11-0640}.

\bibitem[\citeproctext]{ref-Ricketts2018}
Ricketts, Christopher J, Aguirre A De Cubas, Huihui Fan, et al. 2018.
{``The Cancer Genome Atlas Comprehensive Molecular Characterization of
Renal Cell Carcinoma.''} \emph{Cell Reports} 23 (1): 313--326.e5.
\url{https://doi.org/10.1016/j.celrep.2018.03.075}.

\bibitem[\citeproctext]{ref-Sung2021}
Sung, Hyuna, Jacques Ferlay, Rebecca L Siegel, et al. 2021. {``Global
Cancer Statistics 2020: {GLOBOCAN} Estimates of Incidence and Mortality
Worldwide for 36 Cancers in 185 Countries.''} \emph{CA: A Cancer Journal
for Clinicians} 71 (3): 209--49.
\url{https://doi.org/10.3322/caac.21660}.

\bibitem[\citeproctext]{ref-Szklarczyk2023}
Szklarczyk, Damian, Rebecca Kirsch, Mikaela Koutrouli, et al. 2023.
{``The {STRING} Database in 2023: Protein--Protein Association Networks
and Functional Enrichment Analyses for Any of 12,535 Organisms.''}
\emph{Nucleic Acids Research} 51 (D1): D638--46.
\url{https://doi.org/10.1093/nar/gkac1000}.

\bibitem[\citeproctext]{ref-Viechtbauer2010}
Viechtbauer, Wolfgang. 2010. {``Conducting Meta-Analyses in {R} with the
Metafor Package.''} \emph{Journal of Statistical Software} 36 (3):
1--48. \url{https://doi.org/10.18637/jss.v036.i03}.

\bibitem[\citeproctext]{ref-Wu2021}
Wu, Tianzhi, Erqiang Hu, Shuangbin Xu, et al. 2021. {``{clusterProfiler}
4.0: {A} Universal Enrichment Tool for Interpreting Omics Data.''}
\emph{The Innovation} 2 (3): 100141.
\url{https://doi.org/10.1016/j.xinn.2021.100141}.

\end{CSLReferences}

\begin{center}\rule{0.5\linewidth}{0.5pt}\end{center}

\emph{Manuscript prepared for submission to Scientific Reports. Word
count: \textasciitilde7,200 words (body text). All data, analysis
scripts, and output files are available at {[}repository URL{]}.}




\end{document}
